% Ported verbatim from paper/top3_joint_moment_bound.tex lines 622-716.
% Only cross-reference wording was adapted; no mathematics was altered.
Minimal dependencies of rank-one powers have a broader theory; see
Li and Liu~\cite{LiLiu}. The following lemma uses the additional binary
separation hypothesis. Its elementary proof is included.
\begin{lemma}[Moment independence from separating binary observations]\label{binary}
Let $T=\{v_i:i\in J\}$ be a finite set of distinct points in $\R^3$ with
$v_{i1}>0$ and $v_i^TLv_i=\mathcal M>0$, where $L$ is positive definite.
Let $\mathcal B$ be a family of functions $b:J\to\{0,1\}$ containing the
constant-one function and separating every pair of distinct indices.
Suppose, for every real $a$ with $\sum_{i\in J}a_iv_iv_i^T=0$, that
\begin{equation}
 \sum_{i\in J}a_i b(i)v_i=0\quad\hbox{for every }b\in\mathcal B. \label{obs}
\end{equation}
Then the matrices $\{v_iv_i^T:i\in J\}$ are linearly independent.
\end{lemma}
\begin{proof}
Suppose instead that $I\subseteq J$ is a support-minimal dependent family
of these matrices, with relation $a_i\ne0$ for all $i\in I$.
Write $s=|I|$. Since $\dim\operatorname{Sym}_3=6$, minimality gives $s\le7$.
The ellipse equation gives $\sum_{I}a_i=0$, while the constant observation
in \eqref{obs} gives $\sum_I a_iv_i=0$.
With $w_i=(1,v_i)^T\in\R^4$, all three identities combine to give
\[
 \sum_{a_i>0}a_iw_iw_i^T=\sum_{a_i<0}(-a_i)w_iw_i^T.
\]
Both sides are positive semidefinite. The range of a positive weighted sum
of $w_iw_i^T$ is precisely the span of those $w_i$:
its kernel consists of vectors perpendicular to every such $w_i$.
Consequently the two spans coincide, and the affine hull $H$ of all $v_i$
in the circuit satisfies
\[
 1+\dim H=\dim\operatorname{span}\{w_i:i\in I\}
 \le\min(|\{a_i>0\}|,|\{a_i<0\}|)\le3.
\]
This is an elementary positive-Gram range argument, not a new general circuit bound.
Distinct points here cannot be proportional: equal ellipse radius forces
the proportionality factor to be $\pm1$, and the positive first coordinate
excludes minus one. Also an affine line meets this positive-definite ellipsoid
in at most two points, by its strictly quadratic restriction to the line.
If $\dim H\le1$, there are at most two points, whose rank-one matrices
are independent by nonproportionality, a contradiction.
It remains that $\dim H=2$.

If $0\in H$, this is a two-dimensional linear plane. Its rank-one matrices
span a space of dimension at most three, so minimality now gives $s\le4$.
For any observation $b$ nonconstant on $I$, \eqref{obs} and the total vector
relation give nonzero-coefficient linear relations on both nonempty sets
$I_1=\{i:b(i)=1\}$ and $I_0=\{i:b(i)=0\}$.
Nonproportionality forces each set to contain at least three points.
This would give $s\ge6$, a contradiction. If instead every observation in
$\mathcal B$ were constant on $I$, then minimality excludes $s=1$ (a
one-element circuit is impossible because $v_iv_i^T\ne0$), so we may choose
distinct $i,j\in I$; since $\mathcal B$ separates every pair of distinct
indices of $J$, some $b\in\mathcal B$ satisfies $b(i)\ne b(j)$,
contradicting constancy on $I$.

If $0\notin H$, write $H=\{v:c\cdot v=1\}$.
Any three distinct points of $H$ on the ellipsoid are linearly independent:
a linear relation has coefficient sum zero after applying $c$, hence would
make the three points affinely collinear, which the line-intersection fact forbids.
A nonconstant observation would therefore require at least four points
on each of $I_1,I_0$, yielding $s\ge8$, again a contradiction.
Thus every observation is constant on $I$. A one-element circuit is impossible
because $v_iv_i^T\ne0$; separation of any two indices in $I$ is contradicted.
\end{proof}

\begin{theorem}[Every actual positive clique has at most six types]\label{six}
For every actual maximizer, and every positive
clique $J$ of its full support, the map
$a\mapsto\sum_{i\in J}a_iv_iv_i^T$ is injective. In particular $|J|\le6$.
If $\{v_i:i\in J\}$ lies in a linear plane, then $|J|\le3$.
If it lies in an affine plane not containing zero, then $|J|\le5$.
\end{theorem}
\begin{proof}
Use the actual observations $b_k(i)=B_{ki}$, for all rows $k$.
For nonempty $J$, a row $k\in J$ supplies the constant-one observation.
They separate indices: $B$ is symmetric and has distinct full rows.
Lemma \ref{binary} now applies by \eqref{rigid}.
The linear-plane refinement follows from $\dim\operatorname{Sym}_2=3$.
For an affine plane $c\cdot v=1$, the nonzero symmetric matrix
$L-\mathcal M cc^T$ is orthogonal to every $v_iv_i^T$. It is nonzero because $L$
has rank three while $\mathcal M cc^T$ has rank at most one. This leaves dimension at most five.
Empty cliques require no argument.
\end{proof}

\begin{corollary}[Proof of Theorem~\ref{structure-main}]\label{twentyfour}
Every maximizer has at most 24 distinct canonical support points.
\end{corollary}
\begin{proof}
Partition its full support by the signs of coordinates two and three,
assigning a zero coordinate to the positive sign. There are at most four parts.
Within each part, the last two coordinate products are nonnegative and the
first-coordinate product is strictly positive. Each part is a positive clique.
Theorem \ref{six} gives $m\le4\cdot6=24$ on its original support. This proves Theorem~\ref{structure-main}.
\end{proof}
